\chapter{掩码语言建模}
\label{lab:E30_nlp_bert}

\noindent\textbf{配套文件：}\path{notebook/E30_nlp_bert.ipynb}\par
\noindent\textbf{教材关联：}\bookxrefshort{ch:tokenization}、\bookxrefshort{ch:mlm}。

\section*{实验目标}
核对词元、嵌入、注意力与预训练目标之间的契约。

\section*{准备及定位}
先阅读本册实验总则，确认Notebook中的依赖与资产单元。原理计算、库接口迁移和真实模型训练可能具有不同资源要求，应分别记录设备、版本及运行范围。

源码定位可从下列函数开始；应结合前置数据与配置单元阅读。
\begin{itemize}
\item \texttt{my\_wordpiece}
\item \texttt{my\_encode}
\end{itemize}

\section*{操作及观察}
\begin{enumerate}
\item 检查WordPiece分词结果与特殊词元，记录三类嵌入的组合。
\item 逐层核对双向注意力形状，与因果掩码对照可见位置。
\item 检查MLM被选位置及标签，区分预训练目标和下游分类头。
\item 迁移到Transformers时比较配置、输入字段与输出对象。
\end{enumerate}

\section*{结果判读及提交}
提交可见性对照和目标位置表。分类准确率不能直接验证MLM实现；双向编码器不能简单按自回归方式生成下一词元。

提交时附上所执行单元范围、环境与制品身份、原始结果及失败说明。将未运行项目明确列出，不能用Notebook中已有输出替代本次执行证据。

相关方法的原始定义与适用前提见\cite{devlinBERTPretrainingDeep2019}；本实验的运行结果仍须按实际配置单独验证。
